\documentclass[10pt,a4paper]{article}

\usepackage[a4paper,top=17mm,bottom=18mm,left=17mm,right=17mm]{geometry}
\usepackage{fontspec}
\usepackage{xcolor}
\usepackage{tabularx}
\usepackage{array}
\usepackage{enumitem}
\usepackage{tcolorbox}
\usepackage{fancyhdr}
\usepackage{lastpage}
\usepackage{titlesec}
\usepackage{microtype}
\usepackage{multicol}
\usepackage{ragged2e}

\IfFontExistsTF{Arial}{
  \setmainfont{Arial}
  \setsansfont{Arial}
}{
  \setmainfont{TeX Gyre Heros}
  \setsansfont{TeX Gyre Heros}
}

\definecolor{AIEPRed}{HTML}{D62027}
\definecolor{AIEPNavy}{HTML}{102A43}
\definecolor{AIEPInk}{HTML}{1F2933}
\definecolor{AIEPMuted}{HTML}{536171}
\definecolor{AIEPLine}{HTML}{D7DEE8}
\definecolor{AIEPSoft}{HTML}{F6F8FB}
\definecolor{AIEPGold}{HTML}{B08A3C}
\definecolor{AIEPGreen}{HTML}{2F855A}

\pagestyle{fancy}
\fancyhf{}
\renewcommand{\headrulewidth}{0pt}
\fancyfoot[L]{\footnotesize\textcolor{AIEPMuted}{Taller 2 - Ciencia del reciclaje}}
\fancyfoot[R]{\footnotesize\textcolor{AIEPMuted}{GeoGreen Escolar Osorno - \thepage/\pageref{LastPage}}}

\setlength{\parindent}{0pt}
\setlength{\parskip}{5pt}
\hyphenpenalty=10000
\exhyphenpenalty=10000
\sloppy
\setlist[itemize]{leftmargin=*,topsep=2pt,itemsep=2pt}
\setlist[enumerate]{leftmargin=*,topsep=2pt,itemsep=2pt}
\renewcommand{\arraystretch}{1.16}

\titleformat{\section}
  {\Large\bfseries\color{AIEPNavy}}
  {}{0pt}{}
\titleformat{\subsection}
  {\normalsize\bfseries\color{AIEPNavy}}
  {}{0pt}{}

\newcolumntype{Y}{>{\raggedright\arraybackslash}X}

\newtcolorbox{infobox}[2][]{
  colback=white,
  colframe=AIEPLine,
  boxrule=0.45pt,
  arc=1.1mm,
  left=3mm,
  right=3mm,
  top=2.4mm,
  bottom=2.4mm,
  before skip=5pt,
  after skip=5pt,
  title={\scriptsize\bfseries\MakeUppercase{#2}},
  coltitle=AIEPMuted,
  attach title to upper={\par\vspace{1mm}},
  #1
}

\newtcolorbox{keybox}[2][]{
  colback=AIEPSoft,
  colframe=AIEPRed,
  boxrule=0.65pt,
  arc=1.1mm,
  left=3mm,
  right=3mm,
  top=2.4mm,
  bottom=2.4mm,
  before skip=6pt,
  after skip=6pt,
  title={\scriptsize\bfseries\MakeUppercase{#2}},
  coltitle=AIEPMuted,
  attach title to upper={\par\vspace{1mm}},
  #1
}

\newcommand{\eyebrow}[1]{%
  {\small\bfseries\textcolor{AIEPRed}{\MakeUppercase{#1}}}\par
  \vspace{7mm}
  {\color{AIEPRed}\rule{18mm}{0.8pt}}\par
  \vspace{9mm}
}

\newcommand{\blocktitle}[2]{%
  \vspace{2mm}
  {\large\bfseries\color{AIEPNavy}#1}\hfill{\small\bfseries\textcolor{AIEPRed}{#2}}\par
  \vspace{1mm}
}

\newcommand{\fieldbox}[2]{%
  \begin{infobox}[height=#2]{#1}
  \vspace*{\fill}
  \end{infobox}
}

\begin{document}
\RaggedRight

\eyebrow{GeoGreen Escolar Osorno}

{\Huge\bfseries\color{AIEPNavy}Taller 2 - Ciencia del reciclaje\par}
\vspace{3mm}
{\Large\color{AIEPMuted}De basura a material: mirar, separar y decidir mejor\par}

\vspace{9mm}
\begin{tabular}{@{}>{\bfseries\color{AIEPNavy}}l@{\hspace{7mm}}l@{}}
Sesión & Taller 2 \\
Duración & 90 minutos por bloque \\
Modalidad & Presencial, participativa y grupal \\
Organización & Equipos de 5 a 6 estudiantes \\
Producto & Ficha de material conectada con el problema del Taller 1 \\
\end{tabular}

\vspace{8mm}
\begin{keybox}{La gran idea}
Un residuo no es solo basura. Es un material con propiedades, condiciones y posibilidades. Reciclar bien empieza antes de botar: cuando se reconoce el material, se separa, se mantiene limpio y se entrega correctamente.
\end{keybox}

\begin{tabularx}{\textwidth}{@{}Y Y@{}}
\begin{infobox}[height=35mm]{En este taller aprenderás}
\begin{itemize}
  \item Distinguir objeto y material.
  \item Reconocer familias de materiales.
  \item Explicar limpio, seco y separado.
  \item Clasificar residuos con criterio.
\end{itemize}
\end{infobox}
&
\begin{infobox}[height=35mm]{Al finalizar}
\begin{itemize}
  \item Tu equipo tendrá una ficha de material.
  \item Esa ficha se conectará con el problema ambiental trabajado en el Taller 1.
  \item Quedarán preguntas para preparar el Taller 3.
\end{itemize}
\end{infobox}
\end{tabularx}

\newpage
\section{Objetivos de la clase}

\begin{keybox}{Objetivo general}
Al terminar esta clase, podrás explicar qué hace que un material sea reciclable o no: su composición, la separación en origen y la contaminación. También podrás caracterizar un material concreto vinculado al problema ambiental que tu equipo definió en el Taller 1.
\end{keybox}

\subsection{Objetivos específicos}
\begin{enumerate}
  \item Distinguir los grandes tipos de material presentes en residuos cotidianos: plásticos, papel y cartón, vidrio, metales, orgánico y envases multicapa.
  \item Explicar por qué la mezcla y la contaminación reducen o anulan la reciclabilidad de un material.
  \item Aplicar criterios básicos de clasificación con justificación, reconociendo cuándo un material genera duda.
  \item Construir en equipo una ficha que caracterice un material o residuo real observado en el entorno.
  \item Conectar la ficha de material con la frase de problema ambiental formulada en el Taller 1.
\end{enumerate}

\subsection{Competencias transversales}
\begin{tabularx}{\textwidth}{@{}p{37mm}Y@{}}
\textbf{\color{AIEPNavy}Pensamiento científico} & Mirar un residuo como un material con propiedades, no solo como un objeto que se bota. \\
\textbf{\color{AIEPNavy}Pensamiento crítico} & Pasar de ``esto se recicla'' a ``esto se recicla bajo ciertas condiciones''. \\
\textbf{\color{AIEPNavy}Trabajo colaborativo} & Registrar acuerdos, distribuir roles y justificar decisiones de clasificación. \\
\textbf{\color{AIEPNavy}Comunicación efectiva} & Explicar brevemente por qué un material puede recuperarse o perderse. \\
\textbf{\color{AIEPNavy}Responsabilidad ambiental} & Comprender que reciclar mejor depende de hábitos, información y separación oportuna. \\
\end{tabularx}

\newpage
\section{Bloque 1: Nada es solo basura}
\blocktitle{Mirar el residuo por dentro}{20 minutos}

\begin{infobox}{Objetivo del bloque}
Cambiar la forma de mirar un residuo: pasar de verlo como un objeto que se bota a entenderlo como un material con propiedades.
\end{infobox}

\subsection{Objeto y material no son lo mismo}
Cuando botas una botella, no botas solo ``una botella'': botas plástico PET. El objeto es lo que usamos; el material es lo que define si ese residuo se puede recuperar, cuánto vale, cómo se procesa o si termina como basura común.

\begin{keybox}{Idea clave}
Dos objetos distintos pueden estar hechos del mismo material. Y dos objetos parecidos pueden estar hechos de materiales diferentes.
\end{keybox}

\subsection{Familias de materiales}
\begin{tabularx}{\textwidth}{@{}p{25mm}p{44mm}Y@{}}
\textbf{\color{AIEPNavy}Familia} & \textbf{\color{AIEPNavy}Ejemplos} & \textbf{\color{AIEPNavy}Propiedad clave} \\
\hline
Plásticos & Botellas, envases, bolsas, tapas & No todos son iguales ni se reciclan igual. \\
Papel y cartón & Cuadernos, cajas, hojas & Se arruinan si se mojan o ensucian con comida. \\
Vidrio & Botellas, frascos & Puede reciclarse muchas veces sin perder calidad. \\
Metales & Latas, conservas & Son valiosos para reciclar; el aluminio se recupera fácilmente. \\
Orgánico & Restos de comida, cáscaras & No se recicla como envase; puede compostarse. \\
Multicapa & Cajas de leche, envoltorios metalizados & Mezcla materiales pegados y es difícil de separar. \\
\end{tabularx}

\subsection{El truco está en lo que no se ve}
Algunos residuos engañan: un vaso de café puede parecer papel, pero tener una película plástica; una caja de leche puede parecer cartón, pero estar formada por capas de cartón, plástico y aluminio. Por eso no basta con mirar un residuo por encima.

\begin{infobox}{Preguntas guía}
\begin{itemize}
  \item ¿Qué diferencia hay entre el objeto que botamos y el material del que está hecho?
  \item ¿Qué material aparece más en el colegio?
  \item ¿Qué objeto parece de un material, pero en realidad combina varios?
\end{itemize}
\end{infobox}

\newpage
\section{Bloque 2: Reciclable no es lo mismo que reciclado}
\blocktitle{Entender las condiciones}{20 minutos}

\begin{infobox}{Objetivo del bloque}
Comprender que un material no se recicla solo por ser reciclable: depende de cómo se separa, se mantiene y se entrega.
\end{infobox}

\subsection{La diferencia central}
\begin{tabularx}{\textwidth}{@{}Y Y@{}}
\begin{infobox}[height=31mm]{Reciclable}
Es lo que un material podría llegar a ser si entra correctamente a la cadena de recuperación.
\end{infobox}
&
\begin{infobox}[height=31mm]{Reciclado}
Es lo que de verdad ocurrió con el material después de separarlo, retirarlo y procesarlo.
\end{infobox}
\end{tabularx}

\begin{keybox}{Regla fácil de recordar}
Un material reciclable puede perderse si llega sucio, mojado o mezclado. Por eso la condición mínima es: \textbf{limpio, seco y separado}.
\end{keybox}

\subsection{Limpio, seco y separado}
\begin{tabularx}{\textwidth}{@{}p{24mm}Y Y@{}}
\textbf{\color{AIEPNavy}Condición} & \textbf{\color{AIEPNavy}Por qué importa} & \textbf{\color{AIEPNavy}Qué la rompe} \\
\hline
Limpio & Restos de comida o líquidos contaminan el material. & Envases con restos, cajas con grasa. \\
Seco & La humedad arruina sobre todo papel y cartón. & Papel mojado, cartón húmedo. \\
Separado & La mezcla obliga a separar después, y muchas veces ya no se puede. & Botellas, comida y papel en la misma bolsa. \\
\end{tabularx}

\subsection{Colores de referencia en Chile}
\begin{tabularx}{\textwidth}{@{}p{26mm}p{32mm}Y@{}}
\textbf{\color{AIEPNavy}Color} & \textbf{\color{AIEPNavy}Material} & \textbf{\color{AIEPNavy}Ejemplos} \\
\hline
Azul & Papel y cartón & Diarios, cajas, hojas. \\
Amarillo & Plásticos & Botellas, envases, bolsas. \\
Verde & Vidrio & Botellas y frascos. \\
Gris & Metales & Latas de bebida y conservas. \\
Café & Orgánicos & Restos de comida y cáscaras. \\
Rojo & Peligrosos & Pilas, baterías, aceites usados. \\
\end{tabularx}

\subsection{El sistema}
Separar bien tiene sentido porque existe una cadena: separación en origen, punto limpio o retiro, recicladores y gestores, valorización y nuevo material. En Chile, la Ley REP organiza responsabilidades para envases y embalajes, pero la primera decisión sigue ocurriendo antes de botar.

\newpage
\section{Bloque 3: El material bajo la lupa}
\blocktitle{Investigar en equipo}{30 minutos}

\begin{infobox}{Objetivo del bloque}
Elegir un residuo real, caracterizarlo en una ficha de material y practicar una clasificación con criterio.
\end{infobox}

\subsection{Cada equipo, experto de un material}
El equipo elige un residuo concreto: una botella plástica, lata de bebida, caja de leche, envoltorio de snack, hoja usada, vaso desechable, frasco de vidrio o restos de comida. Lo importante es estudiar un material real, no una categoría demasiado amplia.

\begin{tabularx}{\textwidth}{@{}Y Y@{}}
\begin{infobox}[height=42mm]{Roles sugeridos}
\begin{itemize}
  \item Coordinación.
  \item Registro.
  \item Vocería.
  \item Investigación de material.
  \item Investigación de condiciones.
\end{itemize}
\end{infobox}
&
\begin{infobox}[height=42mm]{Buen criterio}
\begin{itemize}
  \item Elegir algo observable.
  \item Preguntar de qué está hecho.
  \item Identificar qué lo arruina.
  \item Conectarlo con el problema del equipo.
\end{itemize}
\end{infobox}
\end{tabularx}

\subsection{Ficha de material}
\begin{keybox}{Campos mínimos}
\begin{itemize}
  \item Material o residuo elegido.
  \item Familia de material.
  \item ¿De qué está hecho realmente?
  \item ¿Es reciclable? ¿Bajo qué condiciones?
  \item ¿Qué lo arruina o lo contamina?
  \item ¿Qué información falta si genera duda?
  \item Relación con el problema del equipo en el Taller 1.
\end{itemize}
\end{keybox}

\subsection{Clasificación con criterio}
\begin{tabularx}{\textwidth}{@{}p{35mm}Y Y@{}}
\textbf{\color{AIEPNavy}Categoría} & \textbf{\color{AIEPNavy}Qué incluye} & \textbf{\color{AIEPNavy}Pista para decidir} \\
\hline
Reciclable & Material recuperable si está limpio y separado. & ¿De qué familia es? ¿Está limpio? \\
No reciclable & Material que hoy no se recupera. & ¿Está demasiado mezclado o es multicapa? \\
Orgánico & Restos de comida y similares. & ¿Se podría compostar? \\
Manejo especial & Pilas, electrónicos, peligrosos. & ¿Puede dañar a personas o al ambiente? \\
Genera duda & Falta información para decidir. & ¿Qué dato falta? \\
\end{tabularx}

\newpage
\section{Bloque 4: Entender el material para imaginar la solución}
\blocktitle{Conectar y proyectar}{20 minutos}

\begin{infobox}{Objetivo del bloque}
Conectar la ficha de material con la frase de problema del Taller 1 y preparar la mirada tecnológica del Taller 3.
\end{infobox}

\subsection{Conectar material y problema}
La pregunta central del cierre es:

\begin{keybox}{Pregunta del equipo}
Ahora que entendemos el material, ¿qué cambia en la forma de mirar nuestro problema?
\end{keybox}

\subsection{Plantilla de conexión}
\begin{tabularx}{\textwidth}{@{}>{\bfseries\color{AIEPNavy}}p{48mm}Y@{}}
Nuestro problema & Escribir la frase de problema del Taller 1. \\
Material en el centro & Identificar el residuo o material clave. \\
Lo que aprendimos & Explicar cómo se comporta, qué lo arruina y bajo qué condiciones se recupera. \\
Qué cambia ahora & Formular una mirada más precisa del problema. \\
\end{tabularx}

\subsection{Puesta en común}
Cada equipo comparte en 30 a 45 segundos:
\begin{enumerate}
  \item Material investigado.
  \item Hallazgo clave de la ficha.
  \item Conexión con el problema ambiental del equipo.
\end{enumerate}

\subsection{Salida reflexiva}
\begin{infobox}{Para cerrar}
\begin{itemize}
  \item ¿Qué aprendimos sobre nuestro material que antes no sabíamos?
  \item ¿Por qué importa entender el material para reciclar bien?
  \item ¿Qué información todavía nos falta?
\end{itemize}
\end{infobox}

\begin{keybox}{Puente al Taller 3}
Con el problema y el material claros, el siguiente paso será imaginar cómo una tecnología puede medir, alertar o visualizar lo que ocurre con los residuos. GeoGreen aparece como ejemplo: un sensor puede ayudar a saber cuánto se llena un contenedor y cuándo actuar.
\end{keybox}

\newpage
\section{Síntesis final}

\begin{tabularx}{\textwidth}{@{}Y Y@{}}
\begin{infobox}[height=34mm]{Aprendizaje 1}
Un residuo no es un objeto cualquiera: es un material con propiedades.
\end{infobox}
&
\begin{infobox}[height=34mm]{Aprendizaje 2}
Que algo sea reciclable no significa que se recicle; depende de que llegue limpio, seco y separado.
\end{infobox}
\end{tabularx}

\begin{infobox}{Aprendizaje 3}
Entender bien un material permite mirar el problema con más precisión y, desde ahí, imaginar una solución que valga la pena.
\end{infobox}

\section{Producto esperado de la sesión}

Al finalizar el taller, deberían quedar:
\begin{itemize}
  \item ficha de material por equipo;
  \item conexión escrita entre la ficha y la frase de problema del Taller 1;
  \item registro de dudas o información faltante;
  \item fotografías o evidencia del trabajo grupal;
  \item una pregunta para continuar hacia el Taller 3.
\end{itemize}

\newpage
\section{Ficha rápida para completar}

\begin{tabularx}{\textwidth}{@{}Y Y@{}}
\fieldbox{Material o residuo elegido}{14mm}
&
\fieldbox{Familia de material}{14mm}
\end{tabularx}

\fieldbox{¿De qué está hecho realmente?}{18mm}
\fieldbox{¿Es reciclable? ¿Bajo qué condiciones?}{18mm}
\fieldbox{¿Qué lo arruina o contamina?}{18mm}
\fieldbox{¿Qué duda queda abierta?}{17mm}
\fieldbox{Relación con el problema del equipo}{20mm}

\end{document}
